\documentclass[11pt,a4paper]{article}

\usepackage[T1]{fontenc}
\usepackage[utf8]{inputenc}
\usepackage[spanish]{babel}
\addto\captionsspanish{\renewcommand{\tablename}{Tabla}}
\usepackage{lmodern}
\usepackage[margin=2.2cm]{geometry}
\usepackage{booktabs,longtable,array}
\usepackage[hidelinks]{hyperref}

\setlength{\parindent}{0pt}
\setlength{\parskip}{0.6em}

\title{GA-TL-02 --- FORM\\Taller de diseño de la arquitectura de microservicios del PFC}
\author{%
  Lombeida Escaleras Bryan Humberto \and
  García Bazurto Kevin Leonel \and
  Manotoa Alava Carlos Alexander \and
  Reinoso Vélez Eduardo David}
\date{Noviembre de 2026}

\begin{document}

\maketitle

\begin{center}
  Universidad Técnica Estatal de Quevedo --- Facultad de Ciencias de la Computación\\
  Carrera de Ingeniería de Software --- Aplicaciones Distribuidas (ISR-701)\\
  Séptimo nivel, paralelo A
\end{center}

\section{Alcance y estado del documento}

El Proyecto Fin de Curso elegido por el equipo es el Sistema de Revisión Técnica Vehicular (RTV) del Centro de Revisión Técnica Vehicular del Municipio de Valencia. Este documento reúne la tabla de servicios del diseño y la tabla de retroalimentación de la revisión docente.

La Tabla~\ref{tab:servicios} procede de los archivos \path{01-dominio-servicios.md} y \path{02-propiedad-datos.md} de la carpeta \path{docs/arquitectura/} del repositorio. Los cuatro servicios son una propuesta del equipo y todavía no tienen la aprobación de los cuatro integrantes. Los endpoints y eventos que aún no están definidos se completarán desde \path{04-comunicaciones-contratos.md}, en esa misma carpeta.

\section{Tabla de servicios}

\begingroup
\footnotesize
\renewcommand{\arraystretch}{1.15}
\begin{longtable}{@{}p{1.9cm}p{3.3cm}p{4.1cm}p{3.1cm}p{2.1cm}@{}}
\caption{Servicios propuestos, datos que poseen y requisito que atienden.}\label{tab:servicios}\\
\toprule
\textbf{Servicio} & \textbf{Responsabilidad} & \textbf{Datos propios} & \textbf{Endpoints o eventos principales} & \textbf{Requisito funcional}\\
\midrule
\endfirsthead
\toprule
\textbf{Servicio} & \textbf{Responsabilidad} & \textbf{Datos propios} & \textbf{Endpoints o eventos principales} & \textbf{Requisito funcional}\\
\midrule
\endhead
\bottomrule
\endfoot
\textbf{MS-01} Turnos y Pagos &
Gestionar el estado administrativo del trámite: turnos, confirmación de pago y autorización de inicio para turnos elegibles. &
Turnos, habilitación, pago, autorizaciones y claves de idempotencia. Base privada \texttt{rtv\_turnos\_pagos}. &
\texttt{POST /turnos/\{turnoId\}/\allowbreak autorizaciones-\allowbreak inspeccion} (contrato E1 v0.2). Resto por definir. &
Sin RF numerado en RF-01 a RF-10. Pendiente de acordar.\\
\addlinespace
\textbf{MS-02} Inspecciones &
Controlar el inicio y el estado de la inspección: solicitar la autorización y registrar el inicio. &
Inspección, estado, registro de inicio y claves de idempotencia. Base privada \texttt{rtv\_inspecciones}. &
\texttt{POST /inspecciones} (contrato E1 v0.2). Resto por definir. &
Sin RF numerado en RF-01 a RF-10. Pendiente de acordar.\\
\addlinespace
\textbf{MS-03} Identidad y Acceso &
Autenticar al personal, gestionar sesiones, roles y permisos, y ofrecer la consulta unificada de auditoría administrativa. &
Cuentas, credenciales protegidas, sesiones, roles, permisos, asignaciones y auditoría de sus propias operaciones. Base privada \texttt{rtv\_identidad\_acceso} (propuesta). &
Por definir. &
RF-01, RF-02, RF-03 y RF-04.\\
\addlinespace
\textbf{MS-04} Ficha y Clasificación Vehicular &
Registrar y consultar la ficha del vehículo, validar sus datos, detectar el cambio de uso y calcular la categoría. &
Fichas e historial, placa y período, tipo, uso, ejes, peso, catálogos, categoría con su fecha, indicación de impronta y auditoría de la ficha. Base privada \texttt{rtv\_ficha\_vehicular} (propuesta). &
Por definir. &
RF-05 a RF-10. Colabora en RF-04.\\
\end{longtable}
\endgroup

Los diez requisitos RF-01 a RF-10 tienen servicio responsable. MS-01 y MS-02 aún no atienden ningún requisito numerado, así que el equipo debe acordar si se amplían los requisitos o se justifica de otro modo su existencia antes de cerrar esta tabla.

\section{Retroalimentación de la revisión de diseño}

La Tabla~\ref{tab:retro} se completa con las observaciones que el docente haga en la revisión de diseño de 10 minutos. No se registran observaciones hipotéticas.

\begin{table}[h]
\centering
\footnotesize
\renewcommand{\arraystretch}{1.2}
\begin{tabular}{@{}p{5.1cm}p{5.1cm}p{4.6cm}@{}}
\toprule
\textbf{Observación} & \textbf{Cambio realizado en el diagrama} & \textbf{Razón si no se aplicó}\\
\midrule
Pendiente de revisión docente & Pendiente & Pendiente\\
\bottomrule
\end{tabular}
\caption{Observaciones recibidas del docente y su tratamiento.}\label{tab:retro}
\end{table}

\section{Declaración de uso de inteligencia artificial}

Eduardo Reinoso utilizó Claude (Anthropic, modelo Sonnet 5.5) para estructurar este documento y ajustar la redacción. Los datos de la Tabla~\ref{tab:servicios} se contrastaron con los archivos 01 y 02 del repositorio antes de incluirlos.

\end{document}
